\begin{python0}
from solutions import *; clear()
\end{python0}
\sectionthree{Find computation}

Suppose you're a bunch of integer values in variables,
say
\begin{center}
\texttt{a}, \texttt{b}, \texttt{c}, \texttt{d}
\end{center}
and you're interested in who has a particular value. If
the target value is in a, you want to print 0. If it's
in b, you want to print 1, etc. In other words, we think of a has having
position 0 and if the target value is found in a, we think of the target
at position 0. I will save the position in the variable \texttt{index}.
The code would look like this:\\
\begin{console}
if (target == a)
{
    index = 0;
}
else if (target == b)
{
    index = 1;
}
else if (target == c)
{
    index = 2;
}
else if (target == d)
{
    index = 3;
}
std::cout << index << '\n';
\end{console}
If the target is not found in a, b, c, d, we will set \texttt{index} to -1:\\
\begin{console}
if (target == a)
{
    index = 0;
}
else if (target == b)
{
    index = 1;
}
else if (target == c)
{
    index = 2;
}
else if (target == d)
{
    index = 3;
}
else
{
    index = -1;
}
std::cout << index << '\n';
\end{console}

\begin{ex}
Complete the program by prompting the user for
\texttt{a}, \texttt{b}, \texttt{c}, \texttt{d} and target and print
\texttt{index}.
\end{ex}
\begin{ex}
Do the same as above but prompt the user for 5
integer values \texttt{a}, \texttt{b}, \texttt{c}, \texttt{d}, \texttt{e}.
\end{ex}
\begin{ex}
Note that the above finds the position of a value by
searching from the first variable to the last (\texttt{a} to \texttt{d}).
What if you search backwards? For instance if the target value is found
in both \texttt{a} and \texttt{d}, your program should print 3 and not 0.
\end{ex}

Make sure you understand and remember this algorithm.

